\subsection{Advertiser-Facing Surfaces}
\label{subsec:advertiser-surfaces}

The components described so far --- matching, re-ranking, engagement
detection, and the adaptive feedback loop --- are all invisible to the
advertiser. They live inside the request path of the chat service and
communicate only with the assistant and the event store. For the system
to function as a product rather than a research prototype, advertisers
need a surface through which they can supply inventory, shape its
targeting, fund its delivery, and inspect its performance. We divide this
surface into two logically distinct layers: a transactional portal for
managing campaign objects, and a read-only analytics layer for
interpreting the events those campaigns generate.

\subsubsection{Advertiser Portal and Campaign Management}
\label{subsubsec:advertiser-portal}

The advertiser portal is a thin REST layer, mounted under the
\texttt{/portal} prefix of the same FastAPI service that serves the chat
endpoint, which exposes CRUD operations over three nested resources:
\emph{advertisers}, \emph{campaigns}, and \emph{ads}. An advertiser is
the billing principal and owns one or more campaigns; a campaign groups
ads under a shared budget and flight window; an ad is the smallest
servable unit and carries the product metadata, targeting constraints,
and creative text that the matching pipeline consumes.

The resource hierarchy is enforced in the URL structure --- campaigns
are reached as
\texttt{/portal/advertisers/\{advertiser\_id\}/campaigns}, ads as
\texttt{/portal/campaigns/\{campaign\_id\}/ads} --- so that ownership is
unambiguous at every write and authorisation reduces to a prefix check.
Each resource supports the usual five operations (\texttt{POST},
\texttt{GET} list, \texttt{GET} by id, \texttt{PATCH}, and soft delete
via a status field), and every mutation is validated against a Pydantic
schema before it reaches the database. Budgets are represented as two
independent scalars, a total budget and a daily budget, both enforced in
the delivery path by the pacing logic described in
Section~\ref{subsubsec:reranking-business-logic}; a campaign whose daily
spend has been exhausted is filtered out of the candidate set for the
remainder of the day rather than being deleted or re-prioritised.

An important design choice is that the portal writes to the same
PostgreSQL tables that the serving path reads from. There is no separate
``staging'' database and no publish step: a newly created ad becomes
eligible for serving as soon as its row is committed and the embedding
worker has populated its vector in pgvector. This keeps the write path
simple and eliminates a class of consistency bugs in which portal state
and serving state drift, at the cost of requiring the embedding worker
to be online for new ads to become servable. We consider this an
acceptable trade because the embedding worker is a small, stateless
service whose failure is easy to detect and whose catch-up on restart is
bounded by the number of ads created during its outage.

\subsubsection{Analytics API and Dashboards}
\label{subsubsec:analytics-api}

The analytics layer is intentionally a read-only projection over the
events table rather than a second source of truth. Every metric the
advertiser sees --- impressions, clicks, engagements, click-through
rate, engagement rate, estimated revenue --- is computed by a SQL
aggregation over the same \texttt{events} rows that the engagement
detector writes in Section~\ref{subsubsec:llm-judge-engagement}. This
discipline is what allows the feedback loop and the reporting surface
to remain consistent by construction: an engagement that reshapes the
next turn's ad selection is the same row that increments the
engagement count on the dashboard, and no reconciliation job is needed
between them.

The API exposes three families of endpoints. A global summary endpoint,
\texttt{GET /analytics/summary?days=N}, returns headline counters and a
ranked list of top-performing ads over a rolling window and is used
both by the internal admin view and by the top-of-dashboard cards on
the advertiser side. A per-advertiser endpoint,
\texttt{GET /portal/advertisers/\{id\}/analytics}, restricts the same
aggregation to ads owned by a single advertiser and is the primary data
source for the advertiser-facing dashboard. A per-campaign endpoint
drills one level further and returns a daily time series suitable for
the line charts on the campaign detail page. All three endpoints accept
a \texttt{days} parameter and translate it into a SQL
\texttt{WHERE timestamp >= NOW() - INTERVAL} clause; none of them
maintain pre-aggregated tables, which we justify by noting that the
events table is indexed on
\texttt{(ad\_id, event\_type, timestamp)} and that the query volume
from the dashboard is orders of magnitude below the write volume from
the serving path.

The React dashboard that consumes these endpoints is deliberately
unambitious in scope. It provides four pages --- an overview with
headline metrics and a time series chart, an advertisers list, a
campaigns list scoped to a selected advertiser, and a campaign detail
view with per-ad breakdowns --- and uses Recharts for all
visualisations. Its visual language is inherited directly from the
chat interface's stylesheet so that the two surfaces read as a single
product rather than two assemblages. The dashboard performs no client-
side aggregation and no caching beyond React Query's default
stale-while-revalidate behaviour; every number on the screen is a
direct reflection of a server-side query over the events table at the
time the page was loaded.